% called by main.tex

\chapter{Architecture and planning}
\label{ch::chapter3}

\section{Introduction}
In the architecture and planning chapter of this OCR application project, we address the key components needed in order to initiate the project. We will define the application requirements, introduce the architecture design, discuss its sprint planning, the work environment and the strategic layout required to convert raw images of any given document type into structured, usable data. \\ 

\section{Analysis and specification of requirements}
This section is all about understanding the project context that we will be developing. It's where  we aim to identify the actors, as well as the functional and non-functional requirements of our system.

\subsection{Identification of actors}
An actor represents any external entity that directly engages with the system. Identifying actors before conducting a comprehensive system analysis is crucial, as they significantly influence how the system interacts with its surrounding environment. In our system, there are currently two actors.

\subsubsection{User} The user can be anyone who needs to extract data from any given document, saving a significant amount of time and effort.
\subsubsection{Admin}
The admin can supervise user activities, as well as manage uploaded images and exported files by visualizing, adding, or deleting them.
 
This identification of actors ensures a clear understanding of how different entities interact with the system, facilitating a more effective system analysis.

\subsection{Functional requirements}

Functional requirements outline the specific actions and behaviors that a system must perform to fulfill its intended purpose. They define the core functionalities and features that users expect from the system.\\\\
\textbf{\large 1 - Image upload }
\begin{itemize}
    \item The system should allow users to upload one or multiple images of documents.
    \item The system should support various image formats (e.g., JPEG, PNG, TIFF).
\end{itemize}

\vspace{4mm}

\textbf{\large 2 - Image Preprocessing  }
\begin{itemize}
    \item The system should preprocess uploaded images to prepare them for more accurate OCR results.
    \item The preprocessing should include tasks such as noise reduction, contrast adjustment, image resizing, grayscale conversion, adaptive thresholding...
\end{itemize}
\vspace{4mm}

\textbf{\large 3 - OCR Scanning  }
\begin{itemize}
    \item The system should accurately recognize and extract text from different uploaded images.
\end{itemize}
\vspace{4mm}

\textbf{\large 4 - Text Organization  }
\begin{itemize}
    \item The system should detect the document type to organize the extracted text in a coherent and appropriate way.
    \item The system should structure the text according to the document type.
\end{itemize}
\vspace{5mm}
\textbf{\large 5 - Human-Machine Interface (HMI) }
\begin{itemize}
    \item The User interface(UI) should display all the data in a user-friendly way. 
    \item The UI should allow users to upload one or multiple images.
    \item The UI should allow users to edit and submit the extracted text to ensure its accuracy.
    \item The UI should allow users to select needed formats and export the data. 
\end{itemize}
\vspace{4mm}

\textbf{\large 6 - Data Submission }
\begin{itemize}
    \item The system should save the edited text to a database upon user submission.
    \item The system should allow the user to get to the data exporting phase upon submission.\\
\end{itemize}
\vspace{4mm}

\textbf{\large 7 - Data Export }
\begin{itemize}
    \item The system should allow users to select one or multiple export formats (Word, CSV, JSON).
    \item The system should export the organized data in the chosen format(s).
\end{itemize}
\vspace{4mm}

\subsection{Non-functional requirements}
Non-functional requirements encompass the standards and limitations that shape the overall performance and user experience of a system. They dictate factors like system speed, reliability, and User interface (UI) design, ensuring that the system operates smoothly and effectively.
\\\\
\textbf{\large 1 - Performance  }
\begin{itemize}
    \item The system should process and extract text from images within a reasonable time frame.
    \item The system should efficiently handle multiple image uploads and large documents.
\end{itemize}
\vspace{4mm}

\textbf{\large 2 - Usability }
\begin{itemize}
    \item The UI should be intuitive and easy to navigate.
    \item Users should be able to upload images, edit text, and export data with minimal training.
\end{itemize}
\vspace{4mm}

\textbf{\large 3 - Reliability }
\begin{itemize}
    \item The OCR functionality should provide accurate text extraction with high reliability.
\end{itemize}
\vspace{4mm}

\textbf{\large 4 - Scalability }
\begin{itemize}
    \item The system should be scalable to accommodate increasing numbers of users and larger datasets.
\end{itemize}
\vspace{4mm}

\textbf{\large 5 - Security }\\
\begin{itemize}
    \item The system should ensure secure storage and transmission of user data.
    \item Access to the system should be controlled and monitored to prevent unauthorized access.
\end{itemize}
\vspace{2mm}

\newpage

\section{Global analysis}
In this section, we take a closer look at the overall system design. We'll explore key elements like the Static diagrams, which are  Use Case Diagram and Class Diagram and dynamic diagrams like Sequence Diagram . Additionally, we'll present the Product Backlog and Sprint Planning. This comprehensive analysis offers a clear understanding of how all components integrate, ensuring that our project is well-structured and aligned with our objectives for successful development.

\subsection{Static diagrams}
A static diagram portrays the fixed characteristics or attributes of a system or its components without showing any temporal progression.
\subsubsection{Global use case diagram}
The use case diagram is a visual tool that shows how users interact with the entire system, highlighting its main features. It helps to quickly understand what users want to achieve and how they interact with the system.

The figure 2.1 represents the global use case diagram .


\begin{figure}[!htpb]
    \centering
    \includegraphics[width=17cm]{figures/GlobalUseCase.png}
    \caption{Global use case diagram.}
    \label{fig:glob_usecase_diag}
\end{figure}

\subsubsection{Global class diagram}
The class diagram is a visual representation of the structure of the system, showing the classes, attributes, and relationships between them. It helps to understand the organization of the system's components and their interactions.

The figure 2.2 represents the global class diagram . 


\begin{figure}[!htpb]
    \centering
    \includegraphics[width= 17cm]{figures/Sprint 3 Class.png}
    \caption{Global class diagram.}
    \label{fig:glob_class_diag}
\end{figure}

\subsection{Dynamic diagrams}
A dynamic diagram illustrates how a system evolves and behaves over time, capturing its operations, actions, and changes.
\subsubsection{Global sequence diagram}
The sequence diagram provides a visual representation of how different components of the system interact over time, illustrating the flow of messages between them. It helps to understand the chronological order of actions and the dependencies between system components.
We have introduced global sequence diagrams from both the user and admin perspectives.
\newpage
The user's global sequence diagram highlights how each action is processed by the system entities, including the user interface, application and machine learning model API, covering actions from image upload to the exporting phase with alternative scenarios such as error handling and conditionals concerning the user's decisions. As shown in the figure 2.3 below.
\\ 

\begin{figure}[!htpb]
    \centering
    \includegraphics[width= 15.5cm]{figures/Sprint 3 Sequence.drawio.png}
    \caption{User global sequence diagram.}
    \label{fig:glob_seq_diag}
\end{figure}
\newpage
The admin's diagram shows how they interact with the system, from logging in to logging out. It covers different scenarios based on the admin's choices, involving key entities like the admin, application, and database. It is showcased in the figure 2.4 below.
\begin{figure}[!htpb]
    \centering
    \includegraphics[height= 19cm]{figures/Global Sequence Admin.png}
    \caption{Admin global sequence diagram.}
    \label{fig:glob_seq_diag}
\end{figure}

\subsection{Product backlog}
This section outlines the Product Backlog the list of functionalities that need to be completed to achieve the project's goals through a table where each line contains a functionality with its ID, user story, priority and complexity.

The table 2.1 represents the global product backlog of the project.
\begin{longtable}{|>{\centering\arraybackslash}m{0.05\linewidth}|>{\centering\arraybackslash}m{0.15\linewidth}|>{\centering\arraybackslash}m{0.45\linewidth}|>{\centering\arraybackslash}m{0.1\linewidth}|>{\centering\arraybackslash}m{0.15\linewidth}|}
\caption{Product backlog.} \label{tab:product_backlog} \\
\hline
\textbf{ID} & \textbf{Functionality} & \textbf{User Story} & \textbf{Priority} & \textbf{Complexity} \\ \hline
\endfirsthead
\hline
\textbf{ID} & \textbf{Functionality} & \textbf{User Story} & \textbf{Priority} & \textbf{Complexity} \\ \hline
\endhead
\hline
\endfoot
\endlastfoot

1 & Image Upload & As a user, I want to be able to upload one or more images of different types like JPEG or PNG so that I can process various documents easily. & High & Medium \\ \hline
2 & Image Preprocessing & As a developer, I want the system to enhance the quality of uploaded images by resizing, reducing noise, adjusting contrast, grayscale conversion so that the user receives more accurate results. & High & High \\ \hline
3 & OCR Integration & As a user, I want the system to read text from images using OCR so that I can extract information without manually typing it. & High & High \\ \hline
4 & Document Type Detection and Text Organization & As a user, I want the system to recognize different types of documents and organize the extracted text accordingly, making it easier for me to read and understand. & High & High \\ \hline
5 & HMI & As a user, I want the interface to be user-friendly and intuitive, allowing me to upload images, view and edit text, and export data without any confusion or difficulty. & Medium & Low \\ \hline
6 & Data Verification & As a user, I want to review and edit the extracted text to ensure it's accurate and error-free before saving it to the system. & High & Medium \\ \hline
7 & Data Submission & As a user, I want to submit the edited text to the system once I'm satisfied with the changes. & High & High \\ \hline
8 & ETL(Extract, Transform, Load) & As a developer, I want the system to perform ETL operations on the extracted text data, to prepare it for the exporting phase. & High & High \\ \hline
9 & Export Data & As a user, I want to export the organized data in formats like Word, CSV, or JSON so that I can use it in other applications. & High & Medium \\ \hline
10 & Optimize Performance & As a user, I want the system to work efficiently and accurately when processing documents so that I have a good experience without delays or errors. & Medium & Medium \\ \hline
11 & Developing the Admin Interface & As an admin, I want to have control over the users' activities in the application by being able to see the lists of uploaded images and exported files, and if needed, I should be able to delete them. & High & High \\ \hline
\end{longtable}

\subsection{Sprint planning}
An essential component of the Agile Methodology is sprint planning. It is about selecting functionalities from the product backlog, estimating their complexity, and breaking them down into manageable tasks. These tasks are grouped into sprints, which are fixed-duration iterations that typically last two to four weeks.



\textbf{\large Sprint 1 (Week 1-4)  }\\\\
\textbf{\t Description}\\
In Sprint 1, we focused on building the core features for handling images. We enabled users to upload images in formats like JPEG and PNG and enhanced image quality using preprocessing techniques. 
We integrated OCR to automate text extraction from images. Additionally, we created a user-friendly interface for uploading images, viewing and editing text, and exporting data.


\textbf{\t Goals}
\begin{enumerate}
  \item Implement image upload functionality;
  \item Develop image preprocessing features;
  \item Integrate OCR into the system;
  \item Design and develop the UI.
\end{enumerate}

\textbf{\t Tasks}
\begin{enumerate}
  \item Implement image upload feature, including file type validation and error handling;
  \item Develop image preprocessing functionalities, such as resizing, noise reduction, and contrast adjustment;
  \item Integrate OCR technology to extract text from uploaded images;
  \item Design all the UI elements.\\
\end{enumerate}


\textbf{\large Sprint 2 (Week 5-8)}\\\\
\textbf{\t Description}\\
During Sprint 2, we developed the document type detection and text organization and ensured that our system recognizes different types of documents and organizes the data coherently. We also enabled users to verify their data to make sure it's error-free and submit it to the system, which then performs ETL to transform it into usable formats suitable for the exporting phase. 


\textbf{\t Goals}
\begin{enumerate}
  \item Develop features for document type detection and text organization;
  \item Implement data verification functionality;
  \item Enable data submission to the system;
  \item Perform ETL operations on the extracted text data.
\end{enumerate}

\textbf{\t Tasks}
\begin{enumerate}
  \item Implement document type detection and text structuring feature;
  \item Develop the verification and edting functionality;
  \item Implement ETL algorithms to prepare data for export.
\end{enumerate}

\textbf{\large Sprint 3 (Week 9-12)}\\\\
\textbf{\t Description}\\
In Sprint 3, we focused on finalizing the project. We enabled users to export the organized data in formats like Word, CSV, and JSON. We optimized the system's performance for efficient and accurate results. we have also built an admin interface to keep the users activities under control.
Finally, we created clear documentation to help users and future developers understand and use the system effectively.

\textbf{\t Goals}
\begin{enumerate}
  \item Implement export functionality for different formats (Word, CSV, JSON);
  \item Optimize system performance for accurate scanning and processing;
  \item Implement an admin interface for the application.
\end{enumerate}
\textbf{\t Tasks}
\begin{enumerate}
  \item Develop export functionality to allow users to export data in Word, CSV, or JSON formats;
  \item Optimize system algorithms and processes to improve performance and accuracy;
  \item Develop the admin interface.
\end{enumerate}

\vspace*{4mm}

\section{Work environment}
In this section, we will outline the different aspects of our work environment that play an important role in the development and success of our project. This includes understanding the modeling language as well as the hardware and software environments that enable us to deliver a high-quality product.

\subsection{Software architectural pattern}
Model-Template-View (MTV) is a software architectural pattern commonly used in web development frameworks like Django. It provides a structured way for designing web applications, dividing the application's logic into three components: the Model, the Template, and the View.\\
The Model represents the data and database interactions, the Template controls how the data is presented visually, and the View manages the interaction between the two, making sure users get the right information in a user-friendly way.\\
The MTV architectural pattern handles how web applications respond to user requests. Here's how it works: When a user visits a specific URL, Django's URL dispatcher (defined in urls.py files) guides the request to the correct view function. The view then interacts with the model, which acts as the intermediary between the application and the database. The model fetches or updates data as needed and returns it to the view. The view then processes this data, selects an appropriate template, and passes the data to it. Using the Django template engine, the template combines the data with pre-existing Hypertext Markup Language (HTML) code to generate the final web page.\\
Finally, the view sends this rendered HTML back to the user as a response.
The figure 2.5 shows the MTV infographic.

\begin{figure}[!htpb]
    \centering
    \includegraphics[height=7cm]{figures/MTV.jpg}
    \caption{MTV infographic.}
    \label{fig:MTV_Infographic}
\end{figure}
\newpage
\subsection{Modeling language}
In our project, we employed the Unified Modeling Language (UML) to visualize and design the system architecture. UML is a widely recognized and standardized modeling language known for its effectiveness in representing complex systems. We selected UML because it enhances communication among team members and involved parties through clear and detailed diagrams. These diagrams help us map out system components, their interactions, and the overall workflow, enabling a detailed understanding of the project's structure. By utilizing UML, we can identify potential issues early in the development process, ultimately contributing to a higher quality final product.
The figure 2.6 shows the UML logo.
\begin{figure}[!htpb]
    \centering
    \includegraphics[height=3cm]{figures/UML_logo.png}
    \caption{UML logo.\cite{UMLLogo}}
    \label{fig:UML_logo}
\end{figure}
\vspace{5mm}
\subsection{Hardware environment}
The hardware environment sets the stage for the performance and reliability of our development activities. Here, we will detail the hardware specifications and configurations that support our development process.

The table 2.2 represents the hardware environment details .

\begin{table}[!htpb]
\caption{Hardware environment table.}
\centering 
\vspace{3 mm}
\def\arraystretch{1.5}
\begin{tabular}{|p{0.2\textwidth}|p{0.3\textwidth}|>{\raggedright\arraybackslash}p{0.3\linewidth}|}
\hline 
\textbf{Brand} & Dell G15 5530&HP Laptop 15-dw3xxx\\ 
\hline 
\textbf{CPU}&  Intel® Core™ i5-13450HX&Intel® Core™ i5-1135G7\\ 
\hline 
\textbf{Installed RAM} &  16GB&12GB\\
\hline 
\textbf{Operating system}  &  Windows 11&Windows 11\\ 
\hline
 \textbf{GPU}& NVIDIA GeForce RTX 3050
&NVIDIA GeForce MX350
\\\hline 
\end{tabular} 
    
    \label{tab:hardware_env_table}
\end{table}
\newpage
\subsection{Software environment}
A strong software environment is essential for effective development and collaboration. This subsection will cover the software tools, frameworks, and programming languages that we have used in this OCR project.


\subsubsection{Development tools}
Development tools are critical in aiding our coding, testing, and debugging efforts. In this subsection, we will describe the specific tools we employ to support our development process.

The table 2.4 shows the different development tools used in this project .
\begin{center}
\vspace{0.5cm}
\captionof{table}{Development tools table.}
\label{tab:Development_tools_table}
\begin{longtable}{|>{\centering\arraybackslash}m{0.3\linewidth}|>{\centering\arraybackslash}m{0.6\linewidth}|}
    \hline
    \textbf{Logo} & \textbf{Description} \\
    \hline
    \endfirsthead
    \hline
    \textbf{Logo} & \textbf{Description} \\
    \hline
    \endhead
    \hline
    \endfoot
    \endlastfoot
    
    \includegraphics[width=3cm]{figures/VSCode_logo.png}&
    \begin{flushleft}
        \textbf{VS Code:} Visual Studio Code is a lightweight but powerful source code editor available for Windows, macOS, and Linux. It comes with built-in support for JavaScript, TypeScript, and Node.js and has a rich ecosystem of extensions for other languages and runtimes.\cite{VSCode} \linebreak \textbf{Usage in Project:} VS Code was utilized as the primary code editor for writing and managing the project's codebase due to its efficiency and extensive features.
    \end{flushleft} 
    \\[2cm] \hline

    \includegraphics[height=3cm]{figures/Git_logo.png} & 
    \begin{flushleft}
        \textbf{Git:} Git is a free and open source distributed version control system designed to handle everything from small to very large projects with speed and efficiency.\cite{Git} \linebreak \textbf{Usage in Project:} Git was employed for version control, ensuring efficient tracking and management of code changes throughout the project's development.
    \end{flushleft} 
    \\[2cm] \hline

    \includegraphics[height=3cm]{figures/Github_logo.png} & 
    \begin{flushleft}
        \textbf{GitHub:} GitHub is a cloud-based platform where you can store, share, and work together with others to write code.\cite{GitHub} \linebreak \textbf{Usage in Project:} GitHub was used in conjunction with Git for version control and collaboration, hosting the project's repository and facilitating teamwork.
    \end{flushleft}
    \\[2cm] \hline
    
    \includegraphics[height=3cm]{figures/Gemini_logo.png} & 
    \begin{flushleft}
        \textbf{Gemini:} Gemini is a family of generative Artificial Intelligence models that lets developers generate content and solve problems. These models are designed and trained to handle both text and images as input.\cite{Gemini} \linebreak \textbf{Usage in Project:} The Gemini Application Programming Interface (API) was used for document type recognition and text organization.
    \end{flushleft} 
    \\[2cm] \hline

    \includegraphics[height=3cm]{figures/Drawio_logo.png } & 
    \begin{flushleft}
        \textbf{Draw.io:} draw.io is a technology stack for building diagramming applications, and the world's most widely used browser-based end-user diagramming software.\cite{DrawioDefinition} \linebreak \textbf{Usage in Project:} Draw.io was used to design the diagrams for the project, helping visualize workflows and system architectures.
    \end{flushleft} 
    \\[2cm] \hline

    \includegraphics[height=3cm]{figures/Illustrator_logo.png} & 
    \begin{flushleft}
        \textbf{Adobe Illustrator:} Adobe Illustrator is a vector graphics editor and design program developed and marketed by Adobe Inc. It is used by graphic designers to create vector images, illustrations, and artwork for various media.\cite{WikipediaIllustrator} \linebreak \textbf{Usage in Project:} Adobe Illustrator was employed to create graphic illustrations for the project, enhancing the visual appeal and clarity of some graphical elements.
    \end{flushleft}
    \\[2cm] \hline

    \includegraphics[height=3cm]{figures/Overleaf_logo.png} & 
    \begin{flushleft}
        \textbf{Overleaf:} Overleaf is an online LaTeX and Rich Text collaborative writing and publishing tool that makes the whole process of writing, editing and publishing scientific documents much quicker and easier.\cite{Overleaf} \linebreak \textbf{Usage in Project:} Overleaf was used to generate this graduation thesis, allowing for collaborative editing and professional typesetting of the document.
    \end{flushleft} 
    \\[2cm] \hline

    \includegraphics[height=3cm]{figures/Notion_logo.png} & 
    \begin{flushleft}
        \textbf{Notion:} Notion is a freemium productivity and note-taking web application developed by Notion Labs, Inc. It offers organizational tools including task management, project tracking, to-do lists, and bookmarking.\cite{Notion}  \linebreak \textbf{Usage in Project:} Notion was used for project management, facilitating team collaboration, tracking project progress and organizing tasks.
    \end{flushleft} 
    \\[2cm] \hline
\end{longtable}
\end{center}



\subsubsection{Libraries, frameworks and programming languages}
Libraries, frameworks and programming languages form the foundation of our software development efforts. This subsection will provide an overview of the frameworks and languages we use.

\vspace{1cm}
\begin{center}
\captionof{table}{Libraries, frameworks and pogramming languages table.}
\label{tab:lib_fram_prog_lang_table}
\begin{longtable}{|>{\centering\arraybackslash}m{0.3\linewidth}|>{\centering\arraybackslash}m{0.6\linewidth}|}
    \hline
    \textbf{Logo} & \textbf{Description} \\
    \hline
    \endfirsthead
    \hline
    \textbf{Logo} & \textbf{Description} \\
    \hline
    \endhead
    \hline
    \endfoot
    \endlastfoot
    
    \includegraphics[height=3cm]{figures/Django_logo.png} & 
    \begin{flushleft}
        \textbf{Django:} Django is a high-level Python web framework that encourages rapid development and clean, pragmatic design.\cite{Django} \linebreak \textbf{Usage in Project:} Django served as the templating engine and framework for the project.
    \end{flushleft} 
    \\[2cm] \hline

    \includegraphics[height=3cm]{figures/Python_logo.png} & 
    \begin{flushleft}
        \textbf{Python:} Python is an interpreted, object-oriented, high-level programming language with dynamic semantics. Its high-level built in data structures, combined with dynamic typing and dynamic binding, make it very attractive for Rapid Application Development, as well as for use as a scripting or glue language to connect existing components together.\cite{Python} \linebreak \textbf{Usage in Project:} Python was the primary programming language used in conjunction with Django to develop the project’s backend.
    \end{flushleft}
    \\[2cm] \hline

    \includegraphics[height=3cm]{figures/html_logo.png} & 
    \begin{flushleft}
        \textbf{HTML:} HTML is the standard markup language for documents designed to be displayed in a web browser. It defines the content and structure of web content. It is often assisted by technologies such as Cascading StyleSheets (CSS) and scripting languages such as JavaScript.\cite{WikipediaHTML} \linebreak \textbf{Usage in Project:} HTML was used to structure the templates in Django, defining the content and layout of the web pages.
    \end{flushleft} 
    \\[2cm] \hline

    \includegraphics[height=3cm]{figures/Css_logo.png} & 
    \begin{flushleft}
        \textbf{CSS:} CSS is a stylesheet language used to describe the presentation of a document written in HTML or XML. CSS describes how elements should be rendered on screen, on paper, in speech, or on other media.\cite{CSS} \linebreak \textbf{Usage in Project:} We used CSS to style the Django templates, ensuring a consistent and visually appealing design.
    \end{flushleft} 
    \\[2cm] \hline

    \includegraphics[height=3cm]{figures/Javascript_logo.png } & 
    \begin{flushleft}
        \textbf{JavaScript:} JavaScript is a lightweight interpreted (or just-in-time compiled) programming language with first-class functions.\cite{JavaScript} \linebreak \textbf{Usage in Project:} JavaScript was utilized to add interactivity and control the behavior of the Django templates.
    \end{flushleft}
    \\[2cm] \hline

    \includegraphics[height=3cm]{figures/Bootstrap_logo.png} & 
    \begin{flushleft}
        \textbf{Bootstrap:} Bootstrap is a free and open-source CSS framework directed at responsive, mobile-first front-end web development. It contains HTML, CSS and JavaScript-based design templates for typography, forms, buttons, navigation, and other interface components.\cite{wikipediaBootstrap}\linebreak \textbf{Usage in Project:} Bootstrap was used to enhance the styling of the Django templates, providing responsive and modern design elements.
    \end{flushleft}
    \\[2cm] \hline

    \includegraphics[height=3cm]{figures/OpenCV_logo.png} & 
    \begin{flushleft}
        \textbf{Open Computer Vision Library (OpenCV):} OpenCV is an open source computer vision and machine learning software library. OpenCV was built to provide a common infrastructure for computer vision applications and to accelerate the use of machine perception in the commercial products.\cite{OpenCV}\linebreak \textbf{Usage in Project:} OpenCV was used for image preprocessing for OCR, preparing images for accurate text recognition.
    \end{flushleft} 
    \\[2cm] \hline

    \includegraphics[width=4cm]{figures/Tesseract_logo.png} & 
    \begin{flushleft}
        \textbf{Tesseract OCR:} Tesseract is an open source OCR platform. OCR extracts text from images and documents without a text layer and outputs the document into a new searchable text file, PDF, or most other popular formats.\cite{Tesseract} \linebreak \textbf{Usage in Project:} Tesseract was employed for the OCR scan, converting images of text into machine-encoded text.
    \end{flushleft} 
    \\[2cm] \hline

\end{longtable}
\end{center}

\section{Conclusion}
In conclusion, the architecture and planning phase laid the foundation for the execution phase of our project. In the next chapters, we'll explore our sprints, where we'll realize this plan through iterative development cycles.
